%% kapa: exam=1 round=21 year=2010 period=1 form=A subject=civil_law subject_ko=민법 seq=13 number=13
%% source: Q-Net 감정평가사 기출문제 (cst00309) artlSeq=5024181 file=extracted/1008/5024181/2010년 제21회 감정평가사 1차시험 기출문제/1교시.hwp
%% conversion: kordoc 4.24.51 -> tools/gen1.py (auto); review=proofread-sonnet; flags=none
\begin{question}{13}
민법상 사단법인의 대표에 관한 설명으로 옳지 않은 것은? (다툼이 있으면 판례에 의함)
\begin{choices}
  \item 법인이 해산한 때에는 파산의 경우를 제외하고는 이사가 청산인이 되는 것이 원칙이다.
  \item 이사가 1인뿐인 법인의 이사가 법인의 재산을 양수하고자 하는 경우, 법원이 선임한 임시이사가 법인을 대표한다.
  \item 정관에 대표권 제한에 관한 규정이 있으나 이를 등기하지 않은 경우, 법인은 그 대표권 제한에 대한 제3자의 선의ㆍ악의에 관계없이 그 제3자에게 대항할 수 없다.
  \item 이사는 정관 또는 총회의 결의로 금지하지 않은 사항에 한하여 타인으로 하여금 특정한 행위를 대리하게 할 수 있다.
  \item 직무대행자는 특별한 사정이 없는 한 법인의 통상 사무에 속하는 행위만을 할 수 있다.
\end{choices}
\end{question}
\answer{2}
